%%%%%%%%%%%%%%%%%%%%%%%%%%%%%%%%%%%%%%%%%%%%%%%%%%%%%%%%%%%%%%%%%%%%%%
% Lớp 4 | Chủ đề 1. Ôn tập và bổ sung | Ôn tập các phép tính trong phạm vi 100 000
% Dạng 4D1?2-3: Điền dấu phép tính; mỗi chữ cái thay một chữ số
%%%%%%%%%%%%%%%%%%%%%%%%%%%%%%%%%%%%%%%%%%%%%%%%%%%%%%%%%%%%%%%%%%%%%%


% nguồn: 2026-11-18_lop4_chu-de-1-tu-luan (nhập 18/11/2026)
\begin{bt}%[4D1N2-3]%[ma:55a46a2c]
% nguon: Toán tư duy lớp 4 – Bài 2 SGK Toán 4 KNTT
	Điền dấu $+$ hoặc $-$ vào mỗi ô trống để được phép tính đúng: $9 \ \square\ 8 \ \square\ 7 = 10$.
	\loigiai{
		Thử các cách điền:\\
		$9 + 8 + 7 = 24$; \quad $9 + 8 - 7 = 10$; \quad $9 - 8 + 7 = 8$; \quad $9 - 8 - 7$ không trừ được.\\
		Vậy phép tính đúng là $9 + 8 - 7 = 10$.
	}
\end{bt}

% nguồn: 2026-11-18_lop4_chu-de-1-tu-luan (nhập 18/11/2026)
\begin{bt}%[4D1H2-3]%[ma:bb2d9631]
% nguon: Toán tư duy lớp 4 – Bài 2 SGK Toán 4 KNTT
	Điền dấu $+$, $-$ hoặc $\times$ vào mỗi ô trống để được phép tính đúng: $6 \ \square\ 6 \ \square\ 6 = 30$.
	\loigiai{
		Kết quả $30$ khá lớn so với $6 + 6 + 6 = 18$ nên cần có phép nhân.\\
		$6 \times 6 = 36$, mà $36 - 6 = 30$.\\
		Vậy phép tính đúng là $6 \times 6 - 6 = 30$.\\
		Lưu ý: trong biểu thức có nhân và trừ, ta làm phép nhân trước.
	}
\end{bt}

% nguồn: 2026-11-18_lop4_chu-de-1-tu-luan (nhập 18/11/2026)
\begin{bt}%[4D1V2-3]%[ma:ee4dee6c]
% nguon: Toán tư duy lớp 4 – Bài 2 SGK Toán 4 KNTT
	Mỗi chữ cái thay cho một chữ số, chữ cái khác nhau thay cho chữ số khác nhau. Tìm $A$ và $B$ biết: $\overline{AB} \times 9 = \overline{A0B}$.
	\loigiai{
		$\overline{AB} \times 9 = (A \times 10 + B) \times 9 = A \times 90 + B \times 9$.\\
		$\overline{A0B} = A \times 100 + B$.\\
		Hai số bằng nhau nên $A \times 90 + B \times 9 = A \times 100 + B$.\\
		Bớt đi phần giống nhau ở hai bên: $B \times 8 = A \times 10$.\\
		$B \times 8$ phải là số tròn chục và $B$ là chữ số nên $B = 5$ (khi đó $5 \times 8 = 40$), suy ra $A = 4$.\\
		Vậy $A = 4$, $B = 5$. Thử lại: $45 \times 9 = 405$.
	}
\end{bt}

% nguồn: 2026-11-18_lop4_chu-de-1-tu-luan-bo-sung (nhập 18/11/2026)
\begin{bt}%[4D1H2-3]%[ma:23bef550]
% nguon: Toán tư duy lớp 4 – Bài 2 SGK Toán 4 KNTT
	Điền dấu $\times$ hoặc $:$ vào mỗi ô trống để được phép tính đúng: $48 \ \square\ 6 \ \square\ 2 = 16$.
	\loigiai{
		Thử các cách điền (tính từ trái sang phải):\\
		$48 : 6 \times 2 = 8 \times 2 = 16$. Đúng!\\
		$48 \times 6 : 2 = 144$; \quad $48 : 6 : 2 = 4$; \quad $48 \times 6 \times 2 = 576$.\\
		Vậy phép tính đúng là $48 : 6 \times 2 = 16$.
	}
\end{bt}
